%! Absolute Value Equations \& Inequalities — Unit 2, Lesson 2.2 — Group Activity (Answer Key)
\documentclass[10pt]{article}
\usepackage{algebra2-article}
\usepackage{algebra2-key}   % pulls in -boxes; do NOT also load -boxes
\newcommand{\TallMath}[1]{$\displaystyle #1\rule[-1.4em]{0pt}{3.2em}$}
\newcommand{\numline}[2]{%
  \draw[->, semithick, charcoal] (#1-0.4,0)--(#2+0.4,0);
  \foreach \x in {#1,...,#2}{\draw[charcoal] (\x,0.10)--(\x,-0.10) node[below, font=\scriptsize]{$\x$};}}

\begin{document}

\pageheader{Unit 2, Lesson 2.2}{Group Activity --- Answer Key}

\vspace{0.08in}

\begin{headlinebox}{forestbg}
\small\textbf{Distance, Cases \& Sets.} An absolute value is a \emph{distance}, so it is never
negative. Solve each problem with your partner, and for every inequality record the answer
\emph{three ways}: set notation, interval notation, and a shaded number line. Isolate the bars
\textbf{before} you split.
\end{headlinebox}

\vspace{0.06in}

\begin{tcolorbox}[colback=white, colframe=black!40, title=\textbf{Tier R --- Remediate},
    fonttitle=\bfseries, arc=1.5mm, left=3mm, right=3mm, top=2mm, bottom=2mm, breakable]
\begin{enumerate}[leftmargin=1.7em, itemsep=8pt]
  \item \textbf{Evaluate.} \quad $|-9| = {\color{keyred}\mathbf{9}}$ \qquad
        $|12| - |-4| = {\color{keyred}\mathbf{8}}$ \qquad $-|{-5}| = {\color{keyred}\mathbf{-5}}$
  \item \textbf{Solve (already isolated).} Each has \emph{two} answers.
        \begin{enumerate}[label=(\alph*), leftmargin=1.8em, itemsep=6pt]
          \item $|x| = 8$: \quad $x = {\color{keyred}\mathbf{8}} \text{ or } {\color{keyred}\mathbf{-8}}$
          \item $|x - 3| = 5$: \quad $x = {\color{keyred}\mathbf{8}} \text{ or } {\color{keyred}\mathbf{-2}}$
          \item $|x + 2| = 6$: \quad $x = {\color{keyred}\mathbf{4}} \text{ or } {\color{keyred}\mathbf{-8}}$
        \end{enumerate}
  \item \textbf{Match.} Draw the two solutions of $|x| = 4$ on the number line below (two dots).
        \begin{center}
        \begin{tikzpicture}[scale=0.5]
          \numline{-6}{6}
          \fill[forest] (-4,0) circle (3.5pt); \fill[forest] (4,0) circle (3.5pt);
        \end{tikzpicture}
        \end{center}
\end{enumerate}
\end{tcolorbox}

\begin{tcolorbox}[colback=white, colframe=black!40, title=\textbf{Tier A --- Approaching Proficiency},
    fonttitle=\bfseries, arc=1.5mm, left=3mm, right=3mm, top=2mm, bottom=2mm, breakable]
\begin{enumerate}[leftmargin=1.7em, itemsep=11pt]
  \item \textbf{Isolate, then solve.} $\;2|x - 1| + 3 = 11$.
        \begin{work}
          2|x - 1| + 3 &= 11 \\
              2|x - 1| &= 8 \\
               |x - 1| &= 4 \\
                 x - 1 &= 4 \quad\text{or}\quad x - 1 = -4 \\
                     x &= 5 \quad\text{or}\quad x = -3
        \end{work}
        Solutions: $x = $ \ans{$5$ or $-3$}
  \item \textbf{``And'' inequality.} $\;|x + 1| < 4$.
        \begin{work}
              |x + 1| &< 4 \\
          -4 < x + 1 &< 4 \\
              -5 < x &< 3
        \end{work}
        Set: \ans{$\{x \mid -5 < x < 3\}$} \qquad Interval: \ans{$(-5,\,3)$}
        \begin{center}
        \begin{tikzpicture}[scale=0.5]
          \numline{-6}{4}
          \draw[very thick, forest] (-5,0)--(3,0);
          \draw[very thick, forest, fill=white] (-5,0) circle (3.5pt);
          \draw[very thick, forest, fill=white] (3,0) circle (3.5pt);
        \end{tikzpicture}
        \end{center}
  \item \textbf{``Or'' inequality.} $\;|2x - 3| \ge 5$.
        \begin{work}
          |2x - 3| &\ge 5 \\
            2x - 3 &\le -5 \quad\text{or}\quad 2x - 3 \ge 5 \\
                2x &\le -2 \quad\text{or}\quad 2x \ge 8 \\
                 x &\le -1 \quad\text{or}\quad x \ge 4
        \end{work}
        Set: \ans{$\{x \mid x \le -1 \text{ or } x \ge 4\}$} \qquad Interval: \ans{$(-\infty,-1]\cup[4,\infty)$}
        \begin{center}
        \begin{tikzpicture}[scale=0.5]
          \numline{-3}{6}
          \draw[very thick, forest, <-] (-3.4,0)--(-1,0);
          \draw[very thick, forest, ->] (4,0)--(6.4,0);
          \fill[forest] (-1,0) circle (3.5pt); \fill[forest] (4,0) circle (3.5pt);
        \end{tikzpicture}
        \end{center}
  \item \textbf{Spot the special case.} One of these has \emph{no solution} and one has \emph{exactly
        one}. Label each and explain in a few words.
        \begin{enumerate}[label=(\alph*), leftmargin=1.8em, itemsep=3pt]
          \item $|x - 2| = -4$: \ans{no solution --- a distance can't be negative}
          \item $|x - 2| = 0$: \ans{exactly one solution, $x=2$ --- distance $0$}
        \end{enumerate}
\end{enumerate}
\end{tcolorbox}

\begin{tcolorbox}[colback=white, colframe=black!40, title=\textbf{Tier E --- Extension},
    fonttitle=\bfseries, arc=1.5mm, left=3mm, right=3mm, top=2mm, bottom=2mm, breakable]
\begin{enumerate}[leftmargin=1.7em, itemsep=11pt]
  \item \textbf{Model a tolerance.} A machine cuts rods to a target length of $20$ cm and accepts any
        rod whose length $L$ is within $0.3$ cm of target.
        \begin{enumerate}[label=(\alph*), leftmargin=1.8em, itemsep=6pt]
          \item Write an absolute-value inequality for an acceptable rod. \ans{$|L - 20| \le 0.3$}
          \item Solve it, and state the \emph{shortest} and \emph{longest} acceptable lengths.
                \begin{work}
                        |L - 20| &\le 0.3 \\
                  -0.3 \le L - 20 &\le 0.3 \\
                      19.7 \le L &\le 20.3
                \end{work}
                Shortest \ans{$19.7$ cm} \quad Longest \ans{$20.3$ cm}
        \end{enumerate}
  \item \textbf{Reason about the constant.} Consider $|x + 5| > k$ and $|x + 5| < k$.
        \begin{enumerate}[label=(\alph*), leftmargin=1.8em, itemsep=6pt]
          \item For which values of $k$ does $|x + 5| < k$ have \emph{no solution}? Explain.
                \ansline{$k \le 0$: a distance is never less than a number that is $0$ or negative.}
          \item For which values of $k$ is \emph{every} real number a solution of $|x + 5| > k$?
                Explain. \ansline{$k < 0$: a distance is always $\ge 0$, which exceeds any negative $k$.}
        \end{enumerate}
\end{enumerate}
\end{tcolorbox}

\end{document}
